\documentclass[11pt,a4paper]{article}
\usepackage{amsmath,amssymb,amsthm}
\usepackage{geometry}
\geometry{margin=1in}
\usepackage{hyperref}
\usepackage{enumitem}

\newtheorem{theorem}{Theorem}[section]
\newtheorem{lemma}[theorem]{Lemma}
\newtheorem{proposition}[theorem]{Proposition}
\newtheorem{corollary}[theorem]{Corollary}
\newtheorem{conjecture}[theorem]{Conjecture}
\theoremstyle{definition}
\newtheorem{definition}[theorem]{Definition}
\newtheorem{problem}[theorem]{Open Problem}
\theoremstyle{remark}
\newtheorem{remark}[theorem]{Remark}

\title{\bfseries The Greedy Dirichlet Sub‑Sum Transform\\
\large Rigorous Foundations, Corrections, and Open Problems}
\author{}
\date{}

\begin{document}
\maketitle

\begin{abstract}
The Greedy Dirichlet Sub‑Sum Transform (GDST) associates to every angle \(\theta\in[0,\pi]\) the Dirichlet series
\[
E(\theta,s)=\sum_{n=0}^{\infty}\frac{\delta_n(\theta)}{(n+2)^s},\qquad 
E_{\pm}(\theta,s)=\sum_{n=0}^{\infty}(-1)^n\frac{\delta_n(\theta)}{(n+2)^s},
\]
where \(\delta_n(\theta)\in\{0,1\}\) are the digits of the greedy harmonic expansion of \(\theta/\pi\).
We correct several inaccuracies in earlier formulations of the programme and build a rigorous foundational theory.
We prove that both series converge conditionally for \(\operatorname{Re}s>0\) and are holomorphic there; we describe the exact set where each digit equals~1 using the greedy harmonic tree; we analyse the regularity of \(E_s(\theta)=E(\theta,s)\) as a function of \(\theta\) and determine its jump structure; we prove the identity linking the signed variant to the Dirichlet eta function; we investigate the possibility of analytic continuation and the growth on vertical lines.
The paper concludes with a list of seven well‑posed open problems that emerge naturally from these results, replacing the earlier invalidated programme.
\end{abstract}

\section{Introduction}
The harmonic–categorical framework of Geere et al.\ \cite{Geere2026a,Geere2026b} introduced a greedy binary expansion of angles using the harmonic fractions \(\pi/(n+2)\).
The \emph{Greedy Dirichlet Sub‑Sum Transform} was subsequently proposed as a one‑parameter family of Dirichlet series that would encode the Riemann zeta function in a spectral problem.
The original programme~\cite{HSTprogramme} claimed that the functions \(E(\theta,s)\) and \(E_{\pm}(\theta,s)\) meromorphically continue to the whole complex plane with poles precisely at the non‑trivial zeros of \(\zeta(s)\).
Those claims rested on an unfinished transfer‑operator construction and on a false monotonicity property of the digit functions \(\delta_n(\theta)\).

In the present paper we rebuild the theory on rigorous, elementary grounds.
We do not reach the Riemann hypothesis; instead we prove a collection of unconditional results that describe the exact analytic structure of the GDST in its natural half‑plane of convergence, correct the earlier mistakes, and set out the genuine open problems that remain.

\section{Preliminaries: the greedy harmonic expansion}
For \(\theta\in[0,\pi]\) set \(x=\theta/\pi\in[0,1]\).
The greedy algorithm with the sequence \(\alpha_n=1/(n+2)\) (\(n\ge 0\)) produces binary digits \(\delta_n(\theta)\in\{0,1\}\) and remainders \(\rho_n\):
\[
\rho_0 = x,\qquad 
\delta_n = \begin{cases}
1 & \text{if }\rho_n \ge \frac{1}{n+2},\\[2pt]
0 & \text{otherwise},
\end{cases}
\qquad 
\rho_{n+1} = \rho_n - \frac{\delta_n}{n+2}.
\]
A simple induction yields the error bound\footnote{For a detailed proof see Programme~I, Lemma~1.}
\begin{equation}\label{eq:err}
0\le \rho_n \le \frac{1}{n+1}\qquad (n\ge 0).
\end{equation}
Consequently the series \(\sum_{n=0}^{\infty}\frac{\delta_n}{n+2}\) converges to \(x\).
The partial sums of the digits
\[
D(N)= \sum_{n=0}^{N-1}\delta_n
\]
satisfy the asymptotic
\begin{equation}\label{eq:Dlog}
D(N) = x\,\log(N) + O(1) \qquad (N\to\infty),
\end{equation}
where the constant in the \(O(1)\) bound can be taken uniform in \(x\) (the proof uses only \eqref{eq:err} and summation by parts).

\section{Definition of the GDST}
\begin{definition}
For \(\theta\in[0,\pi]\) and \(s\in\mathbb{C}\) with \(\operatorname{Re}s>1\) we define
\[
E(\theta,s) \;=\; \sum_{n=0}^{\infty}\frac{\delta_n(\theta)}{(n+2)^s},\qquad
E_{\pm}(\theta,s) \;=\; \sum_{n=0}^{\infty}(-1)^n\frac{\delta_n(\theta)}{(n+2)^s}.
\]
\end{definition}
In the next section we will see that the series converge on the larger half‑plane \(\operatorname{Re}s>0\).

\section{Convergence and analyticity}
\begin{theorem}\label{thm:conv}
For every \(\theta\in[0,\pi]\) the two Dirichlet series \(E(\theta,s)\) and \(E_{\pm}(\theta,s)\) converge (conditionally) for all \(s\) with \(\operatorname{Re}s>0\) and define holomorphic functions on that half‑plane.
\end{theorem}
\begin{proof}
Let \(A(u)=\sum_{0\le n\le u}\delta_n(\theta)\).  From \eqref{eq:Dlog} we have \(|A(u)|\le C_\theta\log(u+2)\).
Applying summation by parts to the partial sum up to \(N\) gives
\[
\sum_{n=0}^{N}\frac{\delta_n}{(n+2)^s}
= \frac{A(N)}{(N+2)^s} + s\int_{0}^{N}\frac{A(t)}{(t+2)^{s+1}}\,dt .
\]
For \(\sigma=\operatorname{Re}s>0\) the boundary term tends to \(0\) because \(|A(N)|\le C\log N\) and \(|(N+2)^{-s}|=N^{-\sigma}\).
The integral converges absolutely as \(N\to\infty\) because \(\int^{\infty}(\log t)\,t^{-\sigma-1}dt<\infty\).
The limit is therefore holomorphic in \(s\) on compact subsets of \(\operatorname{Re}s>0\) by the Weierstrass convergence theorem.
The alternating version is handled identically, using \(|A_{\pm}(u)|=|\sum_{n\le u}(-1)^n\delta_n|\le |A(u)|\le C_\theta\log(u+2)\).
\end{proof}

\section{The set where \(\delta_n=1\) and the greedy harmonic tree}
The digits \(\delta_n(\theta)\) are \emph{not} non‑decreasing functions of \(\theta\); consequently there is no single threshold angle per digit.
The correct description uses the \emph{greedy harmonic tree} from Programme~I.

For a binary prefix \(\varepsilon=(\varepsilon_0,\dots,\varepsilon_{n-1})\in\{0,1\}^n\) define the cylinder interval
\[
I_{\varepsilon} = \{\theta\in[0,\pi] : \delta_k(\theta)=\varepsilon_k\ \text{for } 0\le k<n\}.
\]
These are half‑open intervals; their lengths \(\ell(\varepsilon)\) satisfy a recursive splitting rule:
if \(\ell(\varepsilon)>\alpha_n=\pi/(n+2)\) the node splits, and
\[
I_{\varepsilon 0}=[a_\varepsilon, a_\varepsilon+\ell(\varepsilon)-\alpha_n),\qquad
I_{\varepsilon 1}=[a_\varepsilon+\ell(\varepsilon)-\alpha_n, b_\varepsilon),
\]
whereas if \(\ell(\varepsilon)\le\alpha_n\) the node is unsplit and only the left child exists with the same interval.
All endpoints are finite sums of the harmonic fractions; the set of all split points is dense in \([0,\pi]\).

For a fixed depth \(n\) the condition \(\delta_n(\theta)=1\) means that \(\theta\) lies in the right child of some split node of length \(n\):
\[
\{\theta : \delta_n(\theta)=1\} = \bigsqcup_{\substack{\varepsilon\in\{0,1\}^n\\ \varepsilon\text{ splits}}} I_{\varepsilon 1},
\]
a disjoint union of intervals each of length \(\alpha_n\).

Inserting this into the definition of \(E(\theta,s)\) yields the exact representation
\begin{equation}\label{eq:Etree}
E(\theta,s) = \sum_{n=0}^{\infty} \frac{1}{(n+2)^s}
                \sum_{\substack{\varepsilon\in\{0,1\}^n\\ \varepsilon\text{ splits}}}
                \mathbf{1}_{I_{\varepsilon 1}}(\theta),
\qquad \operatorname{Re}s>0 .
\end{equation}

\section{Regularity in the angle variable}
For a fixed \(s\) write \(E_s(\theta)=E(\theta,s)\).  From \eqref{eq:Etree} we see that \(E_s\) is a limit of step functions.

\begin{theorem}[Jumps of \(E_s\)]\label{thm:jump}
For \(\operatorname{Re}s>1\) the series \eqref{eq:Etree} converges uniformly on \([0,\pi]\); hence \(E_s\) is a regulated function.
Its set of discontinuity points is exactly the set of split points (the endpoints of the right‑child intervals at all depths).
At a split point \(x\) the jump is
\[
\Delta E_s(x) = \sum_{n=0}^{\infty} \frac{1}{(n+2)^s}
                \Bigl( \#\{\varepsilon: |\varepsilon|=n, \; a_{\varepsilon 1}=x\}
                     - \#\{\varepsilon: |\varepsilon|=n, \; b_{\varepsilon 1}=x\} \Bigr).
\]
\end{theorem}
\begin{proof}
Uniform convergence for \(\sigma>1\) follows from the Weierstrass \(M\)-test because the total length of all right‑child intervals at depth \(n\) is at most \(\pi\), so the sum of absolute values of the terms is bounded by \(\sum (n+2)^{-\sigma}\pi<\infty\).
The identification of jumps is immediate from the representation \eqref{eq:Etree} because each indicator contributes a jump of \(+1\) at its left endpoint and \(-1\) at its right endpoint.
\end{proof}

\begin{theorem}[Bounded variation]\label{thm:BV}
For \(\operatorname{Re}s>2\) the function \(E_s\) is of bounded variation on \([0,\pi]\), and its total variation satisfies
\[
V(E_s) \le 2\sum_{n=0}^{\infty} (n+2)^{1-\operatorname{Re}s}.
\]
\end{theorem}
\begin{proof}
The number \(S_n\) of split nodes at depth \(n\) is at most \(n+2\).  Each such node contributes two jumps of size \((n+2)^{-\operatorname{Re}s}\).
Summing the absolute jumps over all split points (ignoring possible cancellations) gives an upper bound for the variation of the jump part.  Since the uniform limit of step functions with uniformly bounded variation has bounded variation, the inequality follows.
\end{proof}

At \(s=1\) the function \(E_1(\theta)=\theta/\pi\) has variation \(1\); for \(1<\operatorname{Re}s\le 2\) the finiteness of the variation is an open question (the bound diverges but the actual variation might be finite due to cancellations).

\section{Relation to the Dirichlet eta function}
The alternating full series is
\[
S_{\text{full}}(s) = \sum_{n=0}^{\infty}\frac{(-1)^n}{(n+2)^s} = 1 - \eta(s),
\]
where \(\eta(s)=(1-2^{1-s})\zeta(s)\) is the Dirichlet eta function, holomorphic for \(\operatorname{Re}s>0\).

Define the complement series
\[
\Omega_{\pm}(\theta,s) = \sum_{n=0}^{\infty} (-1)^n\frac{1-\delta_n(\theta)}{(n+2)^s}.
\]
Its cumulative coefficients are \(A^{\Omega}_N = \sum_{n\le N}(-1)^n - \sum_{n\le N}(-1)^n\delta_n\), bounded by \(1+C_\theta\log(N+2)\); thus \(\Omega_{\pm}\) also converges for \(\operatorname{Re}s>0\).

Pointwise addition of the two convergent series gives:

\begin{theorem}\label{thm:eta}
For every \(\theta\in[0,\pi]\) and \(\operatorname{Re}s>0\),
\[
E_{\pm}(\theta,s) + \Omega_{\pm}(\theta,s) = 1 - \eta(s).
\]
\end{theorem}

This identity partitions the shifted eta series into the selected and rejected parts determined by the greedy expansion.

\section{Analytic continuation}
The generating function
\[
G(\theta,z) = \sum_{n=0}^{\infty} \delta_n(\theta)\, z^{\,n+2},\qquad |z|<1,
\]
satisfies the Mellin relation \(E(\theta,s) = \frac{1}{\Gamma(s)}\int_0^\infty t^{s-1}G(\theta,e^{-t})\,dt\) for \(\operatorname{Re}s>1\).

The coefficients are integers \(0\) or \(1\); by the classical Fatou–Carlson–Polya theorem~\cite{FCP}, a power series with integer coefficients and radius~1 either is rational (coefficients eventually periodic) or has the unit circle as a natural boundary.

\begin{theorem}\label{thm:contin}
For a fixed \(\theta\):
\begin{enumerate}[label=(\roman*)]
\item If the greedy expansion is finite, \(E(\theta,s)\) is a finite Dirichlet polynomial (entire).
\item If the sequence \(\delta_n(\theta)\) is eventually periodic, then \(G(\theta,z)\) is rational and \(E(\theta,s)\) extends to a meromorphic function on \(\mathbb{C}\).
\item Otherwise \(G(\theta,z)\) has the unit circle as a natural boundary, and \(E(\theta,s)\) cannot be meromorphically continued to the whole plane.
\end{enumerate}
\end{theorem}

Whether the line \(\operatorname{Re}s=0\) is a natural boundary for \(E(\theta,s)\) in the third case is a stronger statement; the Mellin transform can sometimes smooth out singularities, but it is highly plausible for almost all \(\theta\) (see Open Problem~\ref{prob:natural}).

\section{Vertical growth}
Using the integral representation
\[
E(\theta,s) = s\int_0^\infty \frac{A(t)}{(t+2)^{s+1}}\,dt,
\qquad A(t)=\sum_{0\le n\le t}\delta_n(\theta),
\]
and the bound \(A(t)= x\log(t+2) + O(1)\), we obtain:

\begin{theorem}
For every \(\theta\) and every \(\sigma_0>0\),
\[
E(\theta,s) = O(|s|) \qquad\text{as } |\operatorname{Im}s|\to\infty,\;
\operatorname{Re}s \ge \sigma_0.
\]
In particular, on the critical line
\[
E(\theta,\tfrac12+it) = O(|t|) \qquad (|t|\to\infty).
\]
The same bounds hold for \(E_{\pm}(\theta,s)\).
\end{theorem}

Combined with the identity \(E_{\pm}+\Omega_{\pm}=1-(1-2^{1-s})\zeta(s)\), this forces a strong cancellation between the two sub‑series on the critical line, since \(\zeta(\frac12+it)\) grows much slower (convexity bound \(O(|t|^{1/6})\)).

\section{Open problems}
The results above lay a rigorous foundation; the following problems emerge naturally and replace the earlier, overly ambitious programme.

\begin{problem}[Classification of periodic sequences]\label{prob:periodic}
Characterise the set of \(\theta\in[0,\pi]\) for which \(\delta_n(\theta)\) is eventually periodic.
Prove that this set is countable and consists precisely of those rational multiples of \(\pi\) whose greedy expansion is finite (or possesses a periodic tail).
Conclude that for almost all \(\theta\) the generating function has the unit circle as a natural boundary.
\end{problem}

\begin{problem}[Natural boundary for the Dirichlet series]\label{prob:natural}
Prove that for almost all \(\theta\) the line \(\operatorname{Re}s=0\) is a natural boundary for \(E(\theta,s)\).
Exhibit a specific \(\theta\) for which this can be demonstrated by showing that the sequence \(\delta_n(\theta)\) contains arbitrarily long squares or other obstructions to analytic continuation.
\end{problem}

\begin{problem}[Meromorphic continuation for exceptional angles]
For those \(\theta\) where the greedy expansion is eventually periodic, determine the meromorphic continuation explicitly.
Are there any such \(\theta\) for which the pole set of \(E(\theta,s)\) includes non‑trivial zeros of \(\zeta(s)\)?
\end{problem}

\begin{problem}[Improved vertical bounds]
Investigate whether \(E(\theta,\frac12+it)=O(|t|^{1/2+\varepsilon})\) or even \(O(|t|^{\varepsilon})\) for some \(\theta\).
Link this to the autocorrelation structure of the digits \(\delta_n\).
\end{problem}

\begin{problem}[Mean‑square asymptotics]
Prove an asymptotic of the form
\[
\int_0^T |E(\theta,\tfrac12+it)|^2\,dt = C_\theta\, T\log T + O(T)
\]
and determine the constant \(C_\theta\) in terms of \(\theta\).
Relate it to the variance of the digit summatory function.
\end{problem}

\begin{problem}[Universality]
Does the family \(\{E(\theta,\frac12+it):t\in\mathbb{R}\}\) possess universality \`a la Voronin?
That is, can it approximate arbitrary analytic functions on compact sets by shifting the imaginary part?
\end{problem}

\begin{problem}[Cancellation structure in the eta identity]
Exploiting the identity \(E_{\pm} + \Omega_{\pm} = 1-\eta\), transfer known bounds for \(\zeta\) (e.g.\ the Lindel\"of hypothesis) into properties of the greedy sub‑sums.
Prove that a strong negative correlation must exist between \(E_{\pm}\) and \(\Omega_{\pm}\) and express it through the geometry of the greedy harmonic tree.
\end{problem}

\vspace{1em}
\noindent\textbf{Acknowledgements.} The author is indebted to Victor Geere for originating the greedy harmonic decomposition and for numerous discussions.  The present work would not exist without the collective effort to correct and solidify the foundations of the harmonic–categorical programme.

\begin{thebibliography}{9}
\bibitem{Geere2026a}
V.~Geere, \emph{A Purely Geometric Reconstruction of the Sine Function}, March 2026.
\url{https://victorgeere.co.za/maths/sine-harmonic.html}
\bibitem{Geere2026b}
V.~Geere, \emph{On the Correlation Structure of a Greedy Harmonic Decomposition of the Sine Function}, March 2026.
\bibitem{Geere2026GHeta}
V.~Geere, \emph{The Greedy Harmonic Decomposition and the Dirichlet Eta Function}, March 2026.
\url{https://victorgeere.co.za/maths/greedy-harmonic-eta.html}
\bibitem{HSTprogramme}
V.~Geere et al., \emph{A Harmonic Sine Transform and the Riemann Zeta Function}, 2026.
\bibitem{FCP}
P.~Fatou, \emph{S\'eries trigonom\'etriques et s\'eries de Taylor}, Acta Math.\ 30 (1906), 335--400.
F.~Carlson, \emph{\"Uber Potenzreihen mit ganzzahligen Koeffizienten}, Math.\ Z.\ 9 (1921), 1--13.
G.~P\'olya, \emph{Untersuchungen \"uber L\"ucken und Singularit\"aten von Potenzreihen}, Math.\ Z.\ 29 (1929), 549--640.
\end{thebibliography}

\end{document}